\documentclass[12pt,twoside]{article}
%\usepackage[table]{xcolor} % important to avoid options clash.
%\input{02465shared_preamble}
%\usepackage{cleveref}
\usepackage{url}
\usepackage{graphics}
\usepackage{multicol}
\usepackage{rotate}
\usepackage{rotating}
\usepackage{booktabs}
\usepackage{hyperref}
\usepackage{pifont}
\usepackage{latexsym}
\usepackage[english]{babel}
\usepackage{epstopdf}
\usepackage{etoolbox}
\usepackage{amsmath}
\usepackage{amssymb}
\usepackage{multirow,epstopdf}
\usepackage{fancyhdr}
\usepackage{booktabs}
\usepackage{xcolor}
\newcommand\redt[1]{ {\textcolor[rgb]{0.60, 0.00, 0.00}{\textbf{ #1} } } }


\newcommand{\m}[1]{\boldsymbol{ #1}}
\newcommand{\EE}{\mathbb{E}}
\DeclareMathOperator*{\argmax}{arg\,max}
\DeclareMathOperator*{\argmin}{arg\,min}
\newcommand{\yoursolution}{ \redt{(your solution here) } } 



\title{ Report 3 hand-in }
\date{ \today }
\author{Alen Hodzic(\texttt{s253790})} 

\begin{document}
\maketitle

\begin{table}[ht!]
\caption{Attribution table.}
\begin{tabular}{llll}
\toprule
                                                                           & Alen Hodzic  &  &  \\
\midrule
 1: Optimal policy                                                         & 100\%  &  &  \\
 2: Simulating a finite approximation of the optimal action-value function & 100\%  &  &  \\
 3: Analytically computing the optimal action-value function               & 100\%  &  &  \\
 4: Extend solution to all states and actions                              & 100\%  &  &  \\
 5: Modified reward function                                               & 100\%  &  &  \\
 6: Long bridge                                                            & 100\%  &  &  \\
\bottomrule
\end{tabular}
\end{table}

%\paragraph{Statement about collaboration:}
%Please edit this section to reflect how you have used external resources. The following statement will in most cases suffice: 
%\emph{The code in the irls/project1 directory is entirely}

%\paragraph{Main report:}
Headings have been inserted in the document for readability. You only have to edit the part which says \yoursolution. 

\section{Jar-Jar at the battle of Naboo (\texttt{jarjar.py})}
\subsubsection*{{\color{red}Problem 3:  Analytically computing the optimal action-value function}}
	
Using the Bellman optimality equations, we obtain
\begin{align}
Q^*(0,1) & = \gamma Q^*(1,-1) \\
Q^*(1,-1) & = -1 + \gamma Q^*(0,1)
\end{align}
therefore,
\begin{align}
Q^*(0,1) & = \frac{-\gamma}{1 - \gamma^2} \\
Q^*(1,-1) & = \frac{-1}{1 - \gamma^2}
\end{align}
	
\subsubsection*{{\color{red}Problem 5:  Modified reward function}}
	
Using the Bellman equation, we can write for both actions:
\begin{align}
v_\pi(s) &= p \big(r_0 + \gamma v_\pi(s+1)\big)
+ (1-p)\big(2r_0 + \gamma v_\pi(s+2)\big), \\
v_\pi(s) &= p \big(r_0 + \gamma v_\pi(s-1)\big)
+ (1-p)\big(2r_0 + \gamma v_\pi(s-2)\big).
\end{align}
Since the reward is independent of the state, the value function does not depend on $s$, i.e. $v_\pi(s)=v$. Therefore,
\begin{align}
v &= p(r_0 + \gamma v) + (1-p)(2r_0 + \gamma v) \\
  &= r_0(2 - p) + \gamma v,
\end{align}
which gives
\begin{align}
v_\pi(s) = \frac{r_0(2 - p)}{1 - \gamma}.
\end{align}

\section{Optimal escape}
\subsubsection*{{\color{red}Problem 6:  Long bridge}}
	
Because the per-step reward satisfies $r_0 < 0$, any unnecessary movement (e.g., moving back and forth) will only decrease the total return. Therefore, the optimal policy will move directly towards one of the exits, either left or right.

If we go left, we need $n$ steps to reach the terminal state. The total return is
\begin{align}
V_{\text{left}} &= r_0 + \gamma r_0 + \gamma^2 r_0 + \cdots + \gamma^{n-1} r_0 + \gamma^n \cdot 1.
\end{align}
Using the formula for the sum of a geometric series, this becomes
\begin{align}
V_{\text{left}} = r_0 \frac{1 - \gamma^n}{1 - \gamma} + \gamma^n.
\end{align}

Similarly, if we go right, we need $m$ steps and obtain reward $2$ at the end:
\begin{align}
V_{\text{right}} &= r_0 + \gamma r_0 + \gamma^2 r_0 + \cdots + \gamma^{m-1} r_0 + \gamma^m \cdot 2 \\
&= r_0 \frac{1 - \gamma^m}{1 - \gamma} + 2\gamma^m.
\end{align}

Therefore, if
\begin{align}
r_0 \frac{1 - \gamma^n}{1 - \gamma} + \gamma^n
>
r_0 \frac{1 - \gamma^m}{1 - \gamma} + 2\gamma^m
\end{align}
one should go left, if
\begin{align}
r_0 \frac{1 - \gamma^n}{1 - \gamma} + \gamma^n
<
r_0 \frac{1 - \gamma^m}{1 - \gamma} + 2\gamma^m
\end{align}
one should go right, and when the two expressions are equal it makes no difference.
\end{document}